\section{Polynomial reduction to one vertex before surface search}
\label{sec:vertex-reduction}

The anchor product in Section~\ref{sec:euler-lp-crushing} need not remain
a global source of exponential work.  The maintained
\path{fast/fastunknot/normal_vertex_reduction.py} constructs a positive
normal ray whenever the current authenticated knot exterior has more than
one vertex.  Repeated verified crushing reaches a one-vertex source, or
already proves that the knot is trivial, before sector discovery begins.
The construction uses no LP or sector enumeration.

The geometric input is Burton--Ozlen's one-vertex construction
\cite[Lemma~6]{burton-ozlen}.  Their expanded edge neighbourhood has a normal
frontier consisting of discs and spheres, including a disc that is not a
vertex link.  The implementation uses this entire bounded-coordinate
frontier, removes its vertex-link components, and applies the existing
positive support extractor.  It need not identify the particular disc used
in the published construction.  This adaptation separates the geometric
existence argument from the independently checked ray and quotient.

\subsection{Local incidence closure and the normal frontier}

The supplied-source validator reconstructs the global vertices, edges and
faces of the finite triangulation.  With several boundary vertices, the
producer chooses a boundary edge joining distinct vertices.  With one
boundary vertex and at least one interior vertex, it chooses an edge from
the boundary vertex to an interior vertex.  The boundary and source
one-skeletons supply these edges by connectedness.  The initial closed
edge neighbourhood meets the boundary in a disc.

Starting from this edge and its endpoints, the producer repeatedly includes
a face whenever at least two of its three local edge occurrences are
already included; inclusion adds all edges and vertices of that face.
Occurrences, rather than distinct global edge identifiers, are essential:
an edge can occur twice around one face of a generalized triangulation.
Glued partner faces share one global face identifier.  After face closure,
every tetrahedron whose four faces are included is absorbed.  Deferring
this last step is equivalent to performing it during closure, since it
introduces no new lower-dimensional cells.

For each unabsorbed tetrahedron, closure permits at most one included face.
That face contributes a normal triangle parallel to it.  If no face is
included, the included edges are absent, a single edge, or two opposite
edges; each contributes its adjacent normal quadrilateral.  Each included
vertex not accounted for by these pieces contributes its vertex triangle.
An included vertex opposite an included face therefore gives a second
parallel triangle.  Two opposite included edges give two parallel copies
of the same quad type, not two incompatible types.

Consequently the constructed vector $x$ has
\[
 0\le x_i\le2,\qquad \sum_i x_i\le4t.
\]
The implementation checks both inequalities, matching, admissibility and
Euler characteristic from the source.  It also records the seed edge,
face additions, absorbed tetrahedra and local piece patterns as diagnostics.

\subsection{Why whole-frontier peeling preserves a positive witness}

Write $V_v$ for the normal vertex link at global vertex $v$, and let $m_v$
be the minimum triangle coordinate of $x$ among its corners.  Set
\[
 x^{\rm can}=x-\sum_v m_vV_v.
\]
This subtraction removes actual disjoint vertex-link components.  A link
can be placed nearest its vertex, disjoint from the remaining normal
surface; the nonnegative matching vector after subtraction realizes the
remainder.  Normal coordinates determine the surface up to normal isotopy,
so the operation does not replace a nonlink component by an unrelated
surface with the same Euler count.

The expanded frontier has only disc and sphere components, and at least
one nonlink disc.  Removing all its vertex-link components leaves that
disc, so
\[
 \chi(x^{\rm can})>0.
\]
There is a triangle zero at every global vertex.  A nonzero triangle-only
matching vector is a nonnegative sum of vertex links; those zeros exclude
such a remainder.  Hence $x^{\rm can}$ also has a quadrilateral.  The
producer checks these conclusions before attempting extraction.

Let $S=\operatorname{supp}(x^{\rm can})$.  The extractor works with every
original matching equation restricted to $S$, using the exact Euler
objective.  Section~\ref{sec:reusable-source-kernels} proves polynomial
support descent to a primitive positive extreme ray $r$, with
$\operatorname{supp}(r)\subseteq S$.  This coordinate face is a face of
the full matching cone.  Thus $r$ is full-source extreme; its primitive
integer coordinates give connectedness.  The inherited zeros preserve
admissibility and exclude a vertex link.  On a classical knot exterior,
the connected positive surface is a disc or a sphere.

\subsection{Progress uses tetrahedra, not the number of vertices}

\begin{proposition}[Complete polynomial vertex reduction]
For a finite triangulation of a classical knot exterior, the maintained
reduction with an unrestricted crushing allowance terminates at a one-vertex
triangulation of the same exterior, or certifies the unknot.  Its bit
complexity is polynomial in the initial explicit triangulation size.
\end{proposition}
\begin{proof}
At a source with several vertices, the preceding construction supplies a
primitive positive nonlink ray.  The crushing theorem of
Section~\ref{sec:normal-crushing-search}, using
\cite[Lemma~5]{burton-ozlen}, applies to any such ray.  A surviving component
with torus boundary is the same knot exterior; if no torus component
survives, the knot is trivial.  At least one quadrilateral is present, so
each nonterminal successor has strictly fewer tetrahedra.  The next source
therefore satisfies the same hypotheses whenever another reduction is
needed.  No step requires the number of vertices itself to decrease.
Strict tetrahedron descent bounds the total number of reductions by the
initial number of tetrahedra.

There are $O(t)$ global faces and at most $O(t)$ successful face additions.
Repeated scans therefore take $O(t^2)$ combinatorial work.  The frontier
and its peeled vector have constant-size coordinates.  Exact source-kernel
preparation, support descent, coordinate validation and quotient replay
have polynomial bit complexity.  The primitive extracted coordinates have
polynomial bit length, and the quotient construction uses their quad
types without expanding their multiplicities.  Each successor is again
an explicit face pairing with fewer tetrahedra.  Summing polynomial work
over the bounded number of successors proves the assertion.
\end{proof}

The published construction has an $O(t^3)$ bound using its own component
extraction procedure.  That exponent is not asserted for this implementation:
our dense exact source elimination and independent replay add polynomial
work whose degree need not match the published algorithm.

\subsection{Integration, certificates and the remaining complexity gap}

The crushing driver enables this preprocessing by default through
\texttt{reduce\_vertices=True}.  It invokes the reduction before sector
search on every current source, including
successors produced by later search rays.  Reduction crushes and search
crushes share the same \texttt{max\_crushes} allowance.  A reached allowance
returns an authenticated partial chain; preprocessing consumes no sector,
corner-candidate, LP-pivot or LP-anchor allowance.  Its work remains part
of the complete running time, and callback cancellation propagates.

There is no new trusted recognition certificate.  Extracted rays use the
existing \texttt{normal-ray-crushing-v1} packet, within the existing
diagram-bound chain and recognition schemas.  The independent checker
validates the source coordinates, primitive full-source extremality,
quotient, retained component and strict descent.  It need not trust the
expansion diagnostics or re-run its discovery procedure.  The low-level
\texttt{ONE\_VERTEX} result reports a vertex count, not a knot verdict;
classical knot-exterior provenance comes from the authenticated chain,
not merely from having torus boundary.

When this preprocessing is enabled, every source passed to sector search
has one vertex.  Thus a supplied sector has at most one active corner
group, and its complete anchor family has at most
$\min(4t,4k+1)$ members for $k$ allowed quad types; an empty active family
still has its single empty tuple.  An envelope call adds at most one LP.
This removes the general Cartesian anchor-product obstruction
unconditionally for these authenticated search sources.  It does not
bound the pivot count of exact Bland simplex, nor compress the complete
$3^t$ quadrilateral sector family.  Those remaining obligations still
prevent a claim of general quasi-polynomial recognition complexity.
